\begin{table*}[tp]
  \centering
  \caption{Adjacent Pair Accuracy by fragment count. Difficulty scales with how many pieces a protein was digested into: the number of possible orderings grows factorially while the evidence available at each junction does not. Lift is the mean paired per-sample difference between the LLM-Guided arm and the Random Order floor within the bin, not a difference of two independent means. Setup: \textnormal{\textnormal{\textit{S. cerevisiae}}}, 20 digestion replicas.}
  \label{tab:stratification_yeast_r20}
  \begin{tabular}{lrrrrrr}
    \toprule
    Fragments & n & Random Order & Fixed Settings & Random Search (no LLM) & LLM-Guided & Lift \\
    \midrule
    2-4 & 1 & 0.000 & 1.000 & 1.000 & 1.000 & +1.000 \\
    5-9 & 17 & 0.100 & 0.735 & 0.844 & 0.861 & +0.761 \\
    10-19 & 19 & 0.122 & 0.541 & 0.673 & 0.696 & +0.574 \\
    20-49 & 43 & 0.046 & 0.440 & 0.563 & 0.593 & +0.547 \\
    50+ & 19 & 0.020 & 0.382 & 0.486 & 0.516 & +0.496 \\
    \bottomrule
  \end{tabular}
  \par\smallskip\footnotesize{n counts proteins in the bin; a bin's mean is taken over those whose ordering metrics are defined (true fragment order recovered), so it can rest on fewer than n.}
\end{table*}
